\documentclass{beamer}

\usepackage[utf8]{inputenc}
\input{../helper}

\title{Introdução à Computação \\ montando uma biblioteca em C}
\author{Alexandre Rademaker}
\date{}

\begin{document}

\begin{frame}
 \maketitle
\end{frame}

\begin{frame}[allowframebreaks,fragile]{montando uma biblioteca}

  Temos uma coleção de funções para manipular listas encadeadas.

  Como transforma-las em uma biblioteca? Como compartilha-las entre
  vários programas que precisem de listas encadeadas?

  Como compilar a biblioteca separadamente dos programas que a
  utilizem?

\end{frame}


\begin{frame}[fragile]{O arquivo list.h}

  \begin{lstlisting}[style=normalc,caption=src/list.h]
    typedef struct node t_node;

    int     list_nth(t_node *list, int n);
    ...
    
    t_node* list_insert(int element, t_node* list);
    t_node* list_reverse(t_node *a);
    t_node* list_create(int arr[], int N);
    t_node* list_append(t_node *a, t_node *b);
    ...

    void    list_nconc(t_node *a, t_node *b);
    void    list_print(t_node *list);
    ...
  \end{lstlisting}

  Neste arquivo vamos \textbf{declarar} os nomes (símbolos).

\end{frame}


\begin{frame}[fragile]{O arquivo list.c}

  \begin{lstlisting}[style=normalc,caption=src/list.c]
    #include <stdio.h>
    #include <stdlib.h>
    #include "list.h"

    typedef struct node {
      int number;
      struct node *next;
    } t_node;

    ...
  \end{lstlisting}

  Neste arquivo vamos \textbf{definir} a estrutura e as funções.

\end{frame}

\begin{frame}[fragile,allowframebreaks]{O arquivo test.c}

  \begin{lstlisting}[style=normalc, caption=src/list-test.c]
    #include <stdio.h>
    #include <stdlib.h>
    #include "list.h"
    
    int main(int argc, char *argv[]) {
      int tmp[] = {10, 20, 30};
      t_node *a = list_create(tmp, 3);
      list_print(a);
      t_node *ac = list_copy(a);
      list_print(ac);
      list_print(list_reverse(ac));
      t_node *s = list_append(a, a); 
      list_print(s);
      list_nconc(ac, ac);
      list_print(ac); // Ops!
    }
  \end{lstlisting}

  \framebreak

  \begin{lstlisting}[style=code,language=lisp]
    CL-USER> (setf *print-circle* t)
    T
    CL-USER> (let ((a (list 1 2 3))
                   (b (list 10 20 30)))
               (values a b (nconc a a)))
    #1=(1 2 3 . #1#)
    (10 20 30)
    #1=(1 2 3 . #1#)
  \end{lstlisting}

  Qual o bug?! Testes que faltam!?
  
\end{frame}

\begin{frame}[fragile,allowframebreaks]{a compilação}

  Imagine que \texttt{list-test.c} é um programa que use nossa
  biblioteca. Poderiam haver outros.

  Queremos separar a compilação da biblioteca da compilação dos programas que usam a biblioteca.

  Afinal, em algum momento nossa biblioteca ficará \textbf{estável} e
  não irá sofrer alterações.

  \framebreak

  Compilando a biblioteca geramos o \texttt{list.o}:

  \texttt{clang -c -Wall list.c}

  Compilando o programa \texttt{list-test.c} que usa a biblioteca e
  incorporando o binário da biblioteca no arquivo binário (executável)
  final.
  
  \texttt{clang -Wall list-test.c list.o -o test}

  \framebreak

  Também podemos não incorporar o binário da biblioteca no executável
  final, criando um \textbf{shared library}, várias outras podem ser
  encontradas no \texttt{/usr/local/lib}.

  \texttt{clang -shared list.o -o list.so}

  seguido de

  \texttt{clang -Wall list-test.c list.so -o test}

  \framebreak

  Mas podemos usar \texttt{make} para nos ajudar!

  O \texttt{make} lê um arquivo \texttt{Makefile} e determina: 1) que
  comandos deve executar; 2) em qual ordem eles devem ser executados.

  Considere o comando:

  \texttt{clang -c -Wall list.c}

  A entrada é o arquivo \texttt{list.c} e ele produz o arquivo
  \texttt{list.o}, então o \texttt{make} saberá que rodar este comando
  só será necessário se a \textbf{data da última alteração} de
  \texttt{list.c} for mais recente que a \textbf{data de criação} de
  \texttt{list.o}.

  \framebreak

  \begin{lstlisting}[style=code,language=make,caption=src/Makefile]
    OPT=-g -Wall ...

    all: test

    list.o: list.c
      clang -c ${OPT} list.c

    list.so: list.o
      clang -shared list.o -o list.so

    test: list.so list.h list-test.c
      clang ${OPT} list-test.c list.so -o test

    clean:
      rm test list.o list.so
  \end{lstlisting}
  
\end{frame}


\end{document}

%%% Local Variables:
%%% mode: latex
%%% TeX-master: t
%%% End:
